\documentclass[10pt]{article}
\usepackage[margin=1in]{geometry}
\usepackage{booktabs}
\usepackage{amsmath}
\usepackage{amssymb}
\begin{document}

\begin{center}
{\large\bfseries ICCAD 2026 \#707 --- Supplementary Results for Author Response}\\[2pt]
\emph{Measured Layer Sensitivity Guides Mixed-Precision ADC Allocation for CIM-Based LLM Inference}
\end{center}

\noindent
These results support the rebuttal only; the paper is unchanged. Unless noted,
the setup is OPT-125M, WikiText-2, p99-clip ADC calibration, paper data split,
INT8 per-channel weights. ``$\Delta$PPL/layer'' is the per-layer measured
sensitivity when a group's ADC bit-width is reduced. We report \emph{denoised,
fixed-seed} measurements; \emph{absolute} $\Delta$PPL magnitudes depend on the
calibration protocol and the number of profiling batches, but the allocation-relevant quantities --- that
$W_{fc2}$ is the most ADC-sensitive group and that the FFN aggregate exceeds the
attention aggregate --- are reproduced. We report group/aggregate quantities; the
single-layer $W_{head}$ group is excluded as numerically unstable (the paper likewise reports the most-sensitive
\emph{non-$W_{head}$} layer, Sec.~5.4).

\vspace{1em}
\noindent\textbf{Table S1. Metric invariance (PPL vs NLL) and weight-quant
robustness} (OPT-125M, INT8 per-channel weights, probe $7{\to}6$).
The FFN aggregate exceeds the attention aggregate, and $W_{fc2}$ is the
most-sensitive group, under \emph{both} PPL- and NLL-based sensitivity. Because
$\mathrm{PPL}=\exp(\overline{\mathrm{NLL}})$ is strictly monotonic with a shared
baseline, the full per-group ranking is provably metric-independent: empirically
the PPL and NLL rankings are identical (Spearman $\rho=1.0$) and the group-level
ILP allocation is unchanged. (Answers \#707C-Q1; supports CR-2 / \#707C-C7 / \#707D-Q1: the
ordering survives a stronger weight-quant floor.)

\begin{center}
\begin{tabular}{lrr}
\toprule
Sensitivity metric & FFN aggregate & Attention aggregate \\
\midrule
$\Delta$PPL/layer (PPL-based) & $+0.070$ & $+0.030$ \\
$\Delta$NLL/layer ($\times10^{-4}$, NLL-based) & $+2.13$ & $+0.93$ \\
\midrule
\multicolumn{3}{l}{$W_{fc2}$ is the single most-sensitive group under both metrics.}\\
\multicolumn{3}{l}{PPL vs NLL ranking: Spearman $\rho=1.00$; group-level ILP allocation identical.}\\
\multicolumn{3}{l}{Values use the fixed paper-split seed; Table~S2 row~1 reports the 3-seed mean.}\\
\bottomrule
\end{tabular}
\end{center}

\vspace{1em}
\noindent\textbf{Table S2. Robustness of the FFN $>$ attention ordering}
(aggregate FFN vs aggregate attention mean $\Delta$PPL/layer; INT8 per-channel
weights). Swept over probe depth and calibration seed on OPT-125M.
FFN $>$ attention holds in 7/9 single-seed OPT settings; seed-averaged it
holds clearly at $7{\to}6$ ($+0.044$ vs $+0.023$) and overwhelmingly at $7{\to}4$
($\approx$11$\times$); with a larger profiling budget (calib 16, eval 60) the
$7{\to}6$ ordering is stable across seeds, indicating the small-probe single-seed
noise is a finite-batch artifact. (Supports CR-1/CR-3; answers \#707A-W1,
\#707C-Q2, \#707D-W3.)

\begin{center}
\begin{tabular}{llrrc}
\toprule
Model & Probe & FFN $\Delta$PPL/l & Attn $\Delta$PPL/l & FFN $>$ Attn \\
\midrule
OPT-125M & $7{\to}6$ (seed mean) & $+0.044$ & $+0.023$ & yes \\
OPT-125M & $7{\to}6$ (3 seeds)   & \multicolumn{2}{c}{2/3 single-seed; 2/2 at calib16/eval60} & 2/3 \\
OPT-125M & $7{\to}5$ (3 seeds)   & \multicolumn{2}{c}{noisy} & 2/3 \\
OPT-125M & $7{\to}4$ (3 seeds)   & $+0.67$ (mean) & $+0.063$ (mean) & 3/3 \\
\midrule
\multicolumn{5}{l}{OPT overall 7/9; clear at $7{\to}6$ seed-mean and $7{\to}4$.}\\
\bottomrule
\end{tabular}
\end{center}

\vspace{1em}
\noindent\begin{minipage}{\textwidth}
\noindent\textbf{Table S3. Full profiling$\to$ILP$\to$PPA pipeline on Qwen2-7B at a
usable operating point} (WikiText-2; \emph{max-clip} ADC calibration --- one of the
paper's three calibration schemes (Table~I) --- 8-bit nominal; baseline PPL
$\approx$16 vs clean $\approx$10). Profiling-guided ILP allocates ADC bits from the
measured group sensitivities under an ADC-area budget; ADC energy follows the same
$2^{b}$ model as in the paper. FFN is far more ADC-sensitive than attention, so the
ILP protects the sensitive group and reduces the insensitive high-column-count
group. Downstream accuracy (PIQA/BoolQ) is preserved within $\approx$1.6\,pp of the
CIM-8b baseline. (Answers \#707C-C8 and \#707A-W2: a complete modern-architecture
pipeline at a realistic operating point.)

\begin{center}
\begin{tabular}{lrrrr}
\toprule
Configuration & ADC area$+$energy & PPL & PIQA & BoolQ \\
\midrule
Clean (fp16, no CIM)                  & ---           & $\approx$10 & 80.4 & 85.2 \\
CIM uniform 8b (baseline)             & ---           & 15.6        & 77.4 & 75.0 \\
\textbf{Profiling-ILP (34\% budget)}  & \textbf{34\%} & \textbf{15.9}\,($+$1.8\%) & \textbf{77.2} & \textbf{73.4} \\
Profiling-ILP (52\% budget)           & 52\%          & 16.4\,($+$4.9\%) & --- & --- \\
\midrule
\multicolumn{5}{l}{Proxy-blind (reduce the \emph{most}-sensitive group): PPL $\to$ 1938 (model destroyed).}\\
\bottomrule
\end{tabular}
\end{center}
\end{minipage}

\vspace{1em}
\noindent\textbf{Analytical mechanism of the inversion (answers \#707A-W1, \#707C-C3/Q2, \#707D-W3).}
A uniform $b$-bit ADC with calibrated full-scale $V_{FS}$ splits its error on a MAC
output $y$ into a \emph{granular} (resolution) term and an \emph{overload/clipping}
term. With saturation rate $s=\Pr[\,|y|>V_{FS}\,]$ and step $\Delta=2V_{FS}/2^{b}$,
\begin{equation*}
D_{\mathrm{ADC}}(b)=\underbrace{D_{\mathrm{gran}}(b)}_{\approx (1-s)\,\Delta^{2}/12\;\propto\;(1-s)\,V_{FS}^{2}\,4^{-b}}+\underbrace{D_{\mathrm{ovl}}}_{\mathbb{E}[(|y|-V_{FS})^{2}\mid |y|>V_{FS}]\;\text{(independent of }b)}.
\end{equation*}
Reducing one bit thus multiplies \emph{only} the granular term by $4$ and leaves
clipping unchanged. To first order the loss impact of reducing group $g$ is
$\;\text{sensitivity}_g\;\approx\;\underbrace{\Delta D_{\mathrm{gran},g}}_{\text{granular budget}}\times\underbrace{J_g}_{\text{residual-stream propagation gain}}$,
where $J_g$ is the gain with which an output perturbation of group $g$ reaches the
loss along the attention/FFN/residual paths (Fig.~5). The saturation rate captures
only the $(1-s)$ part of the granular budget (it ignores $V_{FS}$ and propagation);
the Hessian trace approximates propagation but ignores the budget --- so each
standard proxy sees at most one factor and mispredicts the product, which is
exactly why the paper measures $\Delta$PPL directly.

\vspace{0.8em}
\noindent\textbf{Table S4. Mechanism diagnostics confirm the decomposition}
(OPT-125M, INT8 per-channel weights, p99-clip, group probe $7{\to}6$). This table
isolates the two \emph{mechanistic factors}; the sensitivity ranking itself
(FFN$>$attention; $W_{fc2}$ most sensitive at the paper's operating point) is in
Tables~S1--S2 and paper Table~4. \emph{Decomposition law:} in every group one-bit
reduction scales the granular error $D_{\mathrm{gran}}$ by $4.06\!\approx\!4^{1}$
and the clipping error $D_{\mathrm{ovl}}$ by $1.00$ --- the bit-width touches only
the granular term. \emph{Two anti-correlated factors:} the local granular budget
$\Delta D_{\mathrm{gran}}$ and the propagation gain $J$ are \textbf{perfectly
anti-correlated across groups (Spearman $\rho=-1.0$)}; $W_{qkv}$ receives the
\emph{largest} bit-reduction error yet propagates it $\approx\!400\times$ more
weakly than $W_{fc2}/W_{out}$ ($J$ is hidden-state drift per unit local error;
drift and output KL co-move, $\rho=1.0$) --- the quantitative form of the Fig.~5
attenuation argument and the reason high-saturation attention layers are the
\emph{least} ADC-sensitive. All diagnostics requested by \#707C-C3
(granular/clipping MSE, SQNR, hidden-state drift, output KL) are reported. The
\emph{product} of the two factors reproduces FFN$>$attention; which FFN sub-layer
ranks top is set by this operating-point/architecture-dependent product
($W_{fc2}$ at the paper's max-clip point, $W_{fc1}$ on Qwen2's gated MLP),
explaining the model dependence noted in \#707C-W2.

\begin{center}
\footnotesize
\setlength{\tabcolsep}{5.5pt}
\begin{tabular}{lrrrrrrr}
\toprule
Group & $\kappa$ (kurt.) & $D_{\mathrm{gran}}{\uparrow}$ & $D_{\mathrm{ovl}}{\uparrow}$ & $\Delta D_{\mathrm{gran}}$ & SQNR & KL & $J$ \\
 & heavy tail & ($7{\to}6$) & ($7{\to}6$) & ($\times10^{-4}$) & (dB) & ($\times10^{-3}$) & ($\times10^{3}$) \\
\midrule
$W_{qkv}$ (attn in)  & $1.1$             & $4.06$ & $1.00$ & $1.97$  & $27.5$ & $2.82$ & $0.22$  \\
$W_{out}$ (attn out) & $22$              & $4.06$ & $1.00$ & $0.003$ & $10.8$ & $1.86$ & $127.9$ \\
$W_{fc1}$ (FFN up)   & $47$              & $4.06$ & $1.00$ & $0.256$ & $15.3$ & $4.14$ & $2.09$  \\
$W_{fc2}$ (FFN down) & $1.5\times10^{5}$ & $4.06$ & $1.00$ & $0.004$ & $4.3$  & $1.84$ & $94.2$  \\
\midrule
\multicolumn{8}{l}{Granular budget $\Delta D_{\mathrm{gran}}$ and propagation gain $J$ are anti-correlated across groups (Spearman $\rho=-1.0$).}\\
\multicolumn{8}{l}{Sensitivity $=$ their \emph{product}; neither saturation ($\rho{=}{-}0.70$) nor Hessian ($\rho{=}0.20$) captures both factors.}\\
\bottomrule
\end{tabular}
\end{center}

\end{document}
